% !TeX program = pdflatex
\documentclass[11pt,a4paper]{article}
\usepackage{../../recursos/latex/estilo-notas}

\begin{document}
\cabecalhoaula{07}{Classificação e Classificadores Gaussianos}{AME §7.1--7.3 e §8.1; ISLP §2.2.3 e §4.4}

%==============================================================================
\section*{Objetivos da aula}
%==============================================================================
Ao final desta aula o aluno deve ser capaz de:
\begin{itemize}[leftmargin=1.4cm]
  \item formular o problema de classificação com a \emph{perda 0--1} e derivar o
        \textbf{classificador de Bayes};
  \item descrever a estratégia \emph{plug-in}: estimar $\Prob(Y=c\mid\x)$ e
        classificar pela maior probabilidade;
  \item estimar essa probabilidade pelo Teorema de Bayes: \textbf{Bayes ingênuo}
        (contínuo e discreto), \textbf{LDA} e \textbf{QDA};
  \item entender por que LDA tem fronteira \emph{linear} e QDA \emph{quadrática};
  \item escolher entre LDA e QDA pelo balanço viés--variância.
\end{itemize}

\begin{emsala}
Início do Bloco~II. A boa notícia: \emph{quase tudo} da Parte~I se transfere ---
risco, viés--variância, validação cruzada. A mudança é que $Y$ agora é
\emph{qualitativo} e a perda natural não é a quadrática, mas a 0--1. Pergunta de
abertura: \emph{``qual é o melhor classificador possível, se eu conhecesse a
distribuição dos dados?''}
\end{emsala}

%==============================================================================
\section{O problema de classificação}
%==============================================================================
Agora $Y$ assume valores num conjunto finito de \textbf{classes} $\mathcal{C}$ (ex.:
$\mathcal{C}=\{0,1\}$ no caso binário). Um \textbf{classificador} é uma função
$g:\R^p\to\mathcal{C}$. A perda quadrática não serve: não há $Y-g(\X)$ entre rótulos,
e codificar as classes por números impõe uma ordem que não existe. Só há acerto e
erro, e a perda natural é a \textbf{perda 0--1}, $L(y,g(\x))=\1{y\ne g(\x)}$. O
risco é a \emph{probabilidade de erro}:
\begin{equation}\label{eq:risco01}
  \risco(g) = \E\big[\1{Y\ne g(\X)}\big] = \Prob\big(Y\ne g(\X)\big).
\end{equation}

\begin{emsala}
Antes de escrever a perda 0--1, pergunte: \emph{``por que não usar
$(Y-g(\X))^2$?''} Com três classes codificadas como $1$, $2$ e $3$, confundir a~1
com a~3 custaria quatro vezes mais que confundi-la com a~2. Deixe a turma chegar a
que só existem acerto e erro --- e a perda que registra isso é a indicadora do erro.
\end{emsala}

\begin{teorema}[Classificador de Bayes]\label{teo:bayes}
Sob a perda 0--1, o classificador que minimiza \eqref{eq:risco01} é
\[
  g^*(\x) = \argmax_{c\in\mathcal{C}}\ \Prob(Y=c\mid \X=\x).
\]
No caso binário $\mathcal{C}=\{0,1\}$: $g^*(\x)=1 \iff \Prob(Y=1\mid\x)\ge \tfrac12$.
\end{teorema}
\begin{proof}
(Caso binário, classes $c_1,c_2$.) Condicionando em $\X$,
\[
  \risco(g)=\Prob(Y\ne g(\X))
  =\int \big[\1{g(\x)=c_2}\Prob(Y=c_1\mid\x)+\1{g(\x)=c_1}\Prob(Y=c_2\mid\x)\big]f(\x)\,d\x.
\]
Para \emph{cada} $\x$, o integrando é minimizado escolhendo $g(\x)=c_1$ quando
$\Prob(Y=c_1\mid\x)\ge\Prob(Y=c_2\mid\x)$ e $g(\x)=c_2$ caso contrário --- isto é,
a classe de maior probabilidade a posteriori. Como cada $\x$ é otimizado
pontualmente, o classificador resultante minimiza o risco global.
\end{proof}

\begin{ideia}
O classificador de Bayes $g^*$ é o ``oráculo'': nenhum classificador faz melhor. Seu
risco $R^*$ é o análogo do \emph{erro irredutível} da regressão (Aula~01). Mas
$g^*$ é \emph{desconhecido}, pois depende de $\Prob(Y=c\mid\x)$. Todo método de
classificação é uma tentativa de \emph{aproximá-lo}.
\end{ideia}

\begin{observacao}[O que se transfere da Parte~I]
A estimação do risco \eqref{eq:risco01} é feita por \emph{data splitting}/validação
cruzada (Aula~03), agora medindo a \emph{proporção de erros} na validação. O balanço
viés--variância também vale: $R(g)-R^*$ decompõe-se num termo tipo variância (classe
$\mathcal{G}$ grande demais) e num termo tipo viés (classe pequena demais). E, como
em regressão, compare \emph{poucos} modelos no teste --- faça o \emph{tuning} por CV
dentro do treino (a probabilidade de escolher errado cresce com o número de modelos
comparados, via Hoeffding $+$ união).
\end{observacao}

\begin{atencao}
A acurácia (1 $-$ risco 0--1) pode \textbf{enganar} em dados desbalanceados: com 10
doentes em 1000, o classificador trivial ``todos saudáveis'' acerta 99\% --- e é
inútil. Precisaremos de outras métricas (matriz de confusão, sensibilidade,
precisão) --- tema da Aula~08.
\end{atencao}

%==============================================================================
\section{Classificadores \emph{plug-in}}
%==============================================================================
O Teorema~\ref{teo:bayes} sugere uma receita direta, dita \textbf{\emph{plug-in}}:
\begin{enumerate}[leftmargin=1.6cm,nosep]
  \item estime $\widehat\Prob(Y=c\mid\x)$ para cada classe $c\in\mathcal{C}$;
  \item classifique $\ghat(\x)=\argmax_{c\in\mathcal{C}}\widehat\Prob(Y=c\mid\x)$.
\end{enumerate}
Os métodos diferem em \emph{como} estimam $\Prob(Y=c\mid\x)$ --- a regressão
logística, por exemplo, a modela diretamente. Esta aula desenvolve os que partem do
\textbf{Teorema de Bayes}:
\begin{equation}\label{eq:gen}
  \Prob(Y=c\mid\x) = \frac{f(\x\mid Y=c)\,\Prob(Y=c)}
                          {\sum_{s\in\mathcal{C}} f(\x\mid Y=s)\,\Prob(Y=s)}.
\end{equation}
Estima-se $\Prob(Y=c)$ pelas \emph{proporções amostrais} e $f(\x\mid Y=c)$ por
algum modelo; como o denominador não depende de $c$, basta maximizar o numerador.
Os métodos abaixo se distinguem pelo modelo de $f(\x\mid Y=c)$.

%==============================================================================
\section{Bayes ingênuo}
%==============================================================================
Supõe \textbf{independência condicional} das covariáveis dada a classe:
\begin{equation}\label{eq:nb}
  f(\x\mid Y=c) = \prod_{j=1}^p f(x_j\mid Y=c).
\end{equation}
A suposição quase nunca é literalmente verdadeira, mas é \emph{conveniente} (troca
uma densidade em $\R^p$ por $p$ densidades em $\R$) e frequentemente gera bons
classificadores. Cada $f(x_j\mid Y=c)$ pode ser qualquer modelo univariado, estimado
só com as observações da classe $c$ --- paramétrico ou não. Para covariáveis
contínuas, o mais comum é $X_j\mid Y=c \sim N(\mu_{j,c},\sigma^2_{j,c})$, com média
e variância amostrais por classe e por covariável (\texttt{GaussianNB}). Para
discretas, a fatoração vale com probabilidades, estimadas por proporções dentro de
cada classe: \texttt{BernoulliNB} (binárias), \texttt{CategoricalNB} (categóricas),
\texttt{MultinomialNB} e \texttt{ComplementNB} (contagens, como palavras num texto
--- ponte com a Aula extra E3).

\begin{emsala}
Pergunte: \emph{``que modelo você usaria para $X_j$ se ela fosse um tempo de espera?
Uma contagem? Uma variável binária?''} O ponto é que o Bayes ingênuo não é só
gaussiano: escolher cada $f(x_j\mid Y=c)$ é inferência univariada, e vale olhar os
histogramas por classe antes de escolher.
\end{emsala}

\begin{atencao}
Numericamente, o produto \eqref{eq:nb} de muitos termos pequenos sofre
\emph{underflow} (vira $0$). Trabalhe sempre em \textbf{escala logarítmica}:
$\ghat(\x)=\argmax_c\big[\sum_j \log \widehat f(x_j\mid Y=c)+\log\widehat\Prob(Y=c)\big]$.
\end{atencao}

%==============================================================================
\section{Análise discriminante: LDA e QDA}
%==============================================================================
Aqui modela-se o vetor inteiro como \textbf{gaussiano multivariado}:
$\X\mid Y=c \sim N(\bm{\mu}_c,\bm{\Sigma}_c)$, com densidade
\[
  f(\x\mid Y=c) = \frac{1}{\sqrt{(2\pi)^p\det\bm{\Sigma}_c}}\,
  \exp\Big(-\tfrac12(\x-\bm{\mu}_c)^\top\bm{\Sigma}_c^{-1}(\x-\bm{\mu}_c)\Big).
\]
A distinção:
\begin{itemize}[leftmargin=1.4cm,nosep]
  \item \textbf{QDA} (\emph{Quadratic Discriminant Analysis}): cada classe tem sua
        própria covariância $\bm{\Sigma}_c$.
  \item \textbf{LDA} (\emph{Linear Discriminant Analysis}): supõe covariância
        \emph{comum} $\bm{\Sigma}_c=\bm{\Sigma}$ para todas as classes.
\end{itemize}
Parâmetros por máxima verossimilhança: $\widehat{\bm\mu}_c$ é a média da classe $c$;
no QDA, $\widehat{\bm\Sigma}_c$ é a covariância amostral da classe $c$ (divisor
$\abs{\mathcal{C}_c}$); no LDA, $\widehat{\bm\Sigma}$ reúne as dispersões de todas as
classes em torno das próprias médias, com divisor $n$.

\begin{proposicao}[LDA tem fronteira linear]\label{prop:lda}
Sob $\X\mid Y=c\sim N(\bm{\mu}_c,\bm{\Sigma})$ (covariância comum), o classificador
de Bayes \eqref{eq:gen} equivale a $\ghat(\x)=\argmax_c \delta_c(\x)$, com a
\emph{função discriminante linear}
\[
  \delta_c(\x) = \x^\top\bm{\Sigma}^{-1}\bm{\mu}_c
   - \tfrac12\bm{\mu}_c^\top\bm{\Sigma}^{-1}\bm{\mu}_c + \log\Prob(Y=c).
\]
As fronteiras entre classes $\{\x:\delta_c(\x)=\delta_{c'}(\x)\}$ são
\textbf{hiperplanos}.
\end{proposicao}
\begin{proof}
Classificar por $\argmax_c \Prob(Y=c\mid\x)$ equivale a $\argmax_c \log\big(f(\x\mid
Y=c)\Prob(Y=c)\big)$ (o denominador de \eqref{eq:gen} não depende de $c$). Com a
densidade gaussiana,
$\log f(\x\mid Y=c)=-\tfrac12(\x-\bm{\mu}_c)^\top\bm{\Sigma}^{-1}(\x-\bm{\mu}_c)
+\text{const}$. Expandindo o quadrado, o termo $\x^\top\bm{\Sigma}^{-1}\x$ é
\emph{igual para todas as classes} (pois $\bm{\Sigma}$ é comum) e pode ser
descartado do $\argmax$. Sobram exatamente os termos lineares em $\x$ de
$\delta_c(\x)$. A diferença $\delta_c-\delta_{c'}$ é afim em $\x$, logo cada
fronteira é um hiperplano.
\end{proof}

\begin{observacao}
Em QDA, $\bm{\Sigma}_c$ varia com $c$, então $\x^\top\bm{\Sigma}_c^{-1}\x$ e
$\log\det\bm{\Sigma}_c$ \emph{não} cancelam:
\[
  \delta_c(\x) = -\tfrac12\,\x^\top\bm{\Sigma}_c^{-1}\x
   + \x^\top\bm{\Sigma}_c^{-1}\bm{\mu}_c
   - \tfrac12\bm{\mu}_c^\top\bm{\Sigma}_c^{-1}\bm{\mu}_c
   - \tfrac12\log\det\bm{\Sigma}_c + \log\Prob(Y=c)
\]
é \textbf{quadrática} em $\x$, e as fronteiras são superfícies quádricas.
\end{observacao}

\subsection{LDA ou QDA?}
Com $K$ classes em $\R^p$, o LDA estima $Kp+p(p+1)/2$ parâmetros (médias e uma
covariância); o QDA, $Kp+Kp(p+1)/2$ (uma covariância por classe). Trade-off
viés--variância (Aula~01): QDA é mais flexível (menos viés), mas estima muito mais
parâmetros (mais variância). Prefira LDA com poucas amostras de treino; QDA com
muitas, ou quando a hipótese de covariância comum for inverossímil.

\begin{emsala}
Faça a conta com a turma para $p=10$ e duas classes: $55$ parâmetros de covariância
no LDA, $110$ no QDA. Pergunte: \emph{``com $n=20$, dez observações por classe, o
QDA consegue sequer estimar $\bm{\Sigma}_c$?''} Não: a covariância amostral de dez
pontos em $\R^{10}$ tem posto no máximo $9$ e não é invertível --- e o
\texttt{scikit-learn} recusa o ajuste. A \texttt{Aula prática 07} mostra isso na
tabela do §6.
\end{emsala}

%==============================================================================
\section{Prática em Python}
%==============================================================================
\begin{lstlisting}
from sklearn.discriminant_analysis import (LinearDiscriminantAnalysis as LDA,
                                           QuadraticDiscriminantAnalysis as QDA)
from sklearn.naive_bayes import GaussianNB

for nome, modelo in [("NB", GaussianNB()), ("LDA", LDA()), ("QDA", QDA())]:
    modelo.fit(X_tr, y_tr)
    print(nome, "acuracia:", modelo.score(X_te, y_te))
    # probabilidades plug-in: modelo.predict_proba(X_te)
\end{lstlisting}
O notebook \texttt{Comparação entre classificadores paramétricos} confronta esses
métodos; \texttt{Aula prática 07} calcula o erro de Bayes numa população conhecida
e mede a troca LDA$\times$QDA em $p=2$ e $p=10$.

%==============================================================================
\section*{Resumo}
%==============================================================================
\begin{itemize}[leftmargin=1.4cm]
  \item Perda 0--1 \eqref{eq:risco01}; o \textbf{classificador de Bayes}
        $g^*(\x)=\argmax_c\Prob(Y=c\mid\x)$ é ótimo (Teo.~\ref{teo:bayes}) mas
        desconhecido.
  \item \emph{Plug-in}: estimar $\Prob(Y=c\mid\x)$ e tomar o $\argmax$; aqui, pelo
        Teorema de Bayes \eqref{eq:gen}.
  \item \textbf{Bayes ingênuo}: independência condicional \eqref{eq:nb}, com um
        modelo univariado por covariável, contínua ou discreta.
  \item \textbf{LDA} (covariância comum $\to$ fronteira linear,
        Prop.~\ref{prop:lda}) e \textbf{QDA} (covariância por classe $\to$
        quadrática), com parâmetros por máxima verossimilhança.
  \item LDA $\times$ QDA é viés $\times$ variância: a conta de parâmetros decide
        com poucos dados.
  \item Cuidado com acurácia em dados desbalanceados $\Rightarrow$ Aula~08.
\end{itemize}

\paragraph{Leitura recomendada.} [AME] §7.1 (risco 0--1, classificador de Bayes,
Teorema~9), §7.2--7.3 e §8.1 (\emph{plug-in}, com Bayes ingênuo em §8.1.3 e LDA/QDA
em §8.1.4; as densidades normais de lá omitem o $\tfrac12$ do expoente); [ISLP]
§2.2.3 (classificador de Bayes) e §4.4 (modelos generativos).

\end{document}
